\documentclass[10pt,a4paper]{amsart}
\usepackage{amsmath,amssymb,amsthm,graphicx}
\usepackage{verbatim}
\newenvironment{reference}{\comment}{\endcomment}
%\newenvironment{reference}{}{}
\newenvironment{slogan}{\comment}{\endcomment}
\newenvironment{history}{\comment}{\endcomment}

% For commutative diagrams we use Xy-pic
\usepackage[all]{xy}

% We use 2cell for 2-commutative diagrams.
\xyoption{2cell}
% Theorem environments.
%
\theoremstyle{plain}
\newtheorem{theorem}[subsection]{Theorem}
\newtheorem{proposition}[subsection]{Proposition}
\newtheorem{lemma}[subsection]{Lemma}

\theoremstyle{definition}
\newtheorem{definition}[subsection]{Definition}
\newtheorem{example}[subsection]{Example}
\newtheorem{exercise}[subsection]{Exercise}
\newtheorem{situation}[subsection]{Situation}

\theoremstyle{remark}
\newtheorem{remark}[subsection]{Remark}
\newtheorem{remarks}[subsection]{Remarks}

\numberwithin{equation}{subsection}

% Macros
%
\def\lim{\mathop{\mathrm{lim}}\nolimits}
\def\colim{\mathop{\mathrm{colim}}\nolimits}
\def\Spec{\mathop{\mathrm{Spec}}}
\def\Hom{\mathop{\mathrm{Hom}}\nolimits}
\def\Ext{\mathop{\mathrm{Ext}}\nolimits}
\def\SheafHom{\mathop{\mathcal{H}\!\mathit{om}}\nolimits}
\def\SheafExt{\mathop{\mathcal{E}\!\mathit{xt}}\nolimits}
\def\Sch{\mathit{Sch}}
\def\Mor{\mathop{\mathrm{Mor}}\nolimits}
\def\Ob{\mathop{\mathrm{Ob}}\nolimits}
\def\Sh{\mathop{\mathit{Sh}}\nolimits}
\def\NL{\mathop{N\!L}\nolimits}
\def\CH{\mathop{\mathrm{CH}}\nolimits}
\def\proetale{{pro\text{-}\acute{e}tale}}
\def\etale{{\acute{e}tale}}
\def\QCoh{\mathit{QCoh}}
\def\Ker{\mathop{\mathrm{Ker}}}
\def\Im{\mathop{\mathrm{Im}}}
\def\Coker{\mathop{\mathrm{Coker}}}
\def\Coim{\mathop{\mathrm{Coim}}}

% Boxtimes
%
\DeclareMathSymbol{\boxtimes}{\mathbin}{AMSa}{"02}

%
% Macros for moduli stacks/spaces
%
\def\QCohstack{\mathcal{QC}\!\mathit{oh}}
\def\Cohstack{\mathcal{C}\!\mathit{oh}}
\def\Spacesstack{\mathcal{S}\!\mathit{paces}}
\def\Quotfunctor{\mathrm{Quot}}
\def\Hilbfunctor{\mathrm{Hilb}}
\def\Curvesstack{\mathcal{C}\!\mathit{urves}}
\def\Polarizedstack{\mathcal{P}\!\mathit{olarized}}
\def\Complexesstack{\mathcal{C}\!\mathit{omplexes}}
% \Pic is the operator that assigns to X its picard group, usage \Pic(X)
% \Picardstack_{X/B} denotes the Picard stack of X over B
% \Picardfunctor_{X/B} denotes the Picard functor of X over B
\def\Pic{\mathop{\mathrm{Pic}}\nolimits}
\def\Picardstack{\mathcal{P}\!\mathit{ic}}
\def\Picardfunctor{\mathrm{Pic}}
\def\Deformationcategory{\mathcal{D}\!\mathit{ef}}
\usepackage{cleveref}
\crefname{originallemma}{original lemma}{original lemmas}
\Crefname{originallemma}{Original lemma}{Original lemmas}
\makeatletter
\newlabel{lemma-acyclic-quotient-complexes-bounded-size}{{079K}{source}{}{sourcepointer-1}{}}
\newlabel{lemma-acyclic-quotient-complexes-bounded-size@cref}{{[originallemma][0][]079K}{[1][0][]source}{}{}{}}
\makeatother
\crefname{originalsection}{original section}{original sections}
\Crefname{originalsection}{Original section}{Original sections}
\makeatletter
\newlabel{more-algebra-section-hom-complexes}{{0A8H}{source}{}{sourcepointer-2}{}}
\newlabel{more-algebra-section-hom-complexes@cref}{{[originalsection][0][]0A8H}{[1][0][]source}{}{}{}}
\makeatother
\crefname{originalequation}{original equation}{original equations}
\Crefname{originalequation}{Original equation}{Original equations}
\makeatletter
\newlabel{more-algebra-equation-cohomology-hom-complex}{{0A5X}{source}{}{sourcepointer-3}{}}
\newlabel{more-algebra-equation-cohomology-hom-complex@cref}{{[originalequation][0][]0A5X}{[1][0][]source}{}{}{}}
\makeatother
\crefname{originallemma}{original lemma}{original lemmas}
\Crefname{originallemma}{Original lemma}{Original lemmas}
\makeatletter
\newlabel{homology-lemma-ML-over-ordinals}{{0AAT}{source}{}{sourcepointer-4}{}}
\newlabel{homology-lemma-ML-over-ordinals@cref}{{[originallemma][0][]0AAT}{[1][0][]source}{}{}{}}
\makeatother
\title{A criterion for K-injectivity in a Grothendieck abelian category}
\author{The Stacks Project Authors}
\date{stacks-project:injectives.tex:local-snapshot-20261012; GNU FDL 1.2 or later}
\begin{document}
\maketitle
\noindent\textit{Editorial note.} All delivered mathematical text is the author's own, copied from the pinned Stacks Project chapter file sites-cohomology.tex; no statement or proof step was rewritten, reordered or removed.\par\medskip
\noindent\textit{Editorial note.} The proof is the author's complete argument, each of the successor and limit steps being written out; the only inputs it names are same-chapter and cross-chapter results carried on this card as navigation stubs, and no step is left to the reader.\par\medskip
\noindent\textit{Editorial note.} Every labelled reference occurring inside the counted spans is carried in assembly.external\_references as a non-load-bearing navigation pointer (same-chapter ones at their definition site, cross-chapter ones at the invocation inside the counted span), so no load-bearing mathematics is replaced by a stub. No bibliography citation occurs inside the counted spans.\par\medskip
\noindent\textit{Editorial note.} The same-name check for this card scanned every staged Stacks chapter of this repository -- all sources/stacks-*/pinned-original-source.tex files, the card's own chapter excluded -- and the labels found are recorded in dedup.same\_name\_cross\_chapter on the card.\par\medskip
\noindent\textit{Editorial note.} The preamble is the chapter's own retained minimal preamble (sources/stacks-sites-cohomology/retained-preamble.tex), a verbatim excerpt of the native preamble: the Xy-pic setup and the native theorem declarations with the mathematical macros.\par\medskip
\noindent\textit{Allowed standard prerequisites (editorial).}\begin{itemize}
\item Acyclic complexes of bounded size inside an acyclic complex. Same chapter, Lemma lemma-acyclic-quotient-complexes-bounded-size. Supplies, for a nonzero acyclic quotient, a subcomplex of the cardinal-bounded size that the transfinite construction inserts at successor steps.
\item The complex of morphisms and its cohomology. More on Algebra, Section more-algebra-section-hom-complexes and Equation more-algebra-equation-cohomology-hom-complex. Supply the identification of the complex of morphisms with the groups computing morphisms in the homotopy category and the limit description at a limit ordinal.
\item Mittag-Leffler for ordinals. Homology, Lemma homology-lemma-ML-over-ordinals. Supplies the vanishing of the limit term at a limit ordinal, which closes the transfinite induction.
\item K-injective complexes. Same chapter, the definition of a K-injective complex. Supplies the property that the proof establishes, namely the vanishing of morphisms from every acyclic complex.
\end{itemize}\medskip
\ifdefined\hypertarget\hypertarget{sourcepointer-1}{}\fi
\noindent\textit{Source reference.} Original lemma 079K: sources/stacks-injectives/pinned-original-source.tex, source lines 1808--1808. Navigation pointer invoked at proof line 1853; the counted span contains the pointer only and the cited result is not reproduced inside it, so the pointer is registered as a non-load-bearing navigation stub in the established form for Stacks material. No bibliography citation is involved.\par\medskip
\ifdefined\hypertarget\hypertarget{sourcepointer-2}{}\fi
\noindent\textit{Source reference.} Original section 0A8H: sources/stacks-more-algebra/pinned-original-source.tex, source lines 17851--17851. Navigation pointer invoked at proof line 1909; the counted span contains the pointer only and the cited result is not reproduced inside it, so the pointer is registered as a non-load-bearing navigation stub in the established form for Stacks material. No bibliography citation is involved.\par\medskip
\ifdefined\hypertarget\hypertarget{sourcepointer-3}{}\fi
\noindent\textit{Source reference.} Original equation 0A5X: sources/stacks-more-algebra/pinned-original-source.tex, source lines 17875--17875. Navigation pointer invoked at proof line 1912; the counted span contains the pointer only and the cited result is not reproduced inside it, so the pointer is registered as a non-load-bearing navigation stub in the established form for Stacks material. No bibliography citation is involved.\par\medskip
\ifdefined\hypertarget\hypertarget{sourcepointer-4}{}\fi
\noindent\textit{Source reference.} Original lemma 0AAT: sources/stacks-homology/pinned-original-source.tex, source lines 8079--8079. Navigation pointer invoked at proof line 1917; the counted span contains the pointer only and the cited result is not reproduced inside it, so the pointer is registered as a non-load-bearing navigation stub in the established form for Stacks material. No bibliography citation is involved.\par\medskip
\begin{lemma}
\label{lemma-characterize-K-injective}
Let $\mathcal{A}$ be a Grothendieck abelian category.
Let $\kappa$ be a cardinal as in
Lemma \ref{lemma-acyclic-quotient-complexes-bounded-size}.
Suppose that $I^\bullet$ is a complex such that
\begin{enumerate}
\item each $I^j$ is injective, and
\item for every bounded above acyclic complex $M^\bullet$
such that $|M^n| \leq \kappa$
we have $\Hom_{K(\mathcal{A})}(M^\bullet, I^\bullet) = 0$.
\end{enumerate}
Then $I^\bullet$ is an $K$-injective complex.
\end{lemma}

\begin{proof}
Let $M^\bullet$ be an acyclic complex. We are going to construct by
induction on the ordinal $\alpha$ an acyclic subcomplex
$K_\alpha^\bullet \subset M^\bullet$ as follows.
For $\alpha = 0$ we set $K_0^\bullet = 0$. For $\alpha > 0$
we proceed as follows:
\begin{enumerate}
\item If $\alpha = \beta + 1$ and $K_\beta^\bullet = M^\bullet$
then we choose $K_\alpha^\bullet = K_\beta^\bullet$.
\item If $\alpha = \beta + 1$ and $K_\beta^\bullet \not = M^\bullet$
then $M^\bullet/K_\beta^\bullet$ is a nonzero acyclic complex.
We choose a subcomplex $N_\alpha^\bullet \subset M^\bullet/K_\beta^\bullet$
as in Lemma \ref{lemma-acyclic-quotient-complexes-bounded-size}.
Finally, we let $K_\alpha^\bullet \subset M^\bullet$
be the inverse image of $N_\alpha^\bullet$.
\item If $\alpha$ is a limit ordinal we set
$K_\beta^\bullet = \colim K_\alpha^\bullet$.
\end{enumerate}
It is clear that $M^\bullet = K_\alpha^\bullet$ for a suitably large
ordinal $\alpha$. We will prove that
$$
\Hom_{K(\mathcal{A})}(K_\alpha^\bullet, I^\bullet)
$$
is zero by transfinite induction on $\alpha$. It holds for $\alpha = 0$
since $K_0^\bullet$ is zero. Suppose it holds for $\beta$ and
$\alpha = \beta + 1$. In case (1) of the list above the result is clear.
In case (2) there is a short exact sequence of complexes
$$
0 \to K_\beta^\bullet \to K_\alpha^\bullet \to N_\alpha^\bullet \to 0
$$
Since each component of $I^\bullet$ is injective we see that we obtain
an exact sequence
$$
\Hom_{K(\mathcal{A})}(K_\beta^\bullet, I^\bullet) \leftarrow
\Hom_{K(\mathcal{A})}(K_\alpha^\bullet, I^\bullet) \leftarrow
\Hom_{K(\mathcal{A})}(N_\alpha^\bullet, I^\bullet)
$$
By induction the term on the left is zero and by assumption on $I^\bullet$
the term on the right is zero. Thus the middle group is zero too.
Finally, suppose that $\alpha$ is a limit ordinal. Then we see that
$$
\Hom^\bullet(K_\alpha^\bullet, I^\bullet) =
\lim_{\beta < \alpha} \Hom^\bullet(K_\beta^\bullet, I^\bullet)
$$
with notation as in
More on Algebra, Section \ref{more-algebra-section-hom-complexes}.
These complexes compute morphisms in $K(\mathcal{A})$ by
More on Algebra, Equation
(\ref{more-algebra-equation-cohomology-hom-complex}).
Note that the transition maps in the system are surjective
because $I^j$ is injective for each $j$. Moreover, for a limit
ordinal $\alpha$ we have equality of limit and value
(see displayed formula above). Thus we may apply
Homology, Lemma \ref{homology-lemma-ML-over-ordinals}
to conclude.
\end{proof}

\end{document}
